% =====================================================================
%  Supplementary Note (\input by supplement.tex after the generated UMLS crosswalk):
%  credibility tiers, why UMLS has no proxy, and the external alignment.
%  Generated from the credibility sweep (protocol data/umls/sweeps/
%  credibility.yaml; results are not redistributed -- regenerate with
%  gem-umls-sweep + gem-umls-classify-credibility against a licensed UMLS
%  copy) and data/umls/adjudications.yaml (credibility/*). Review before use.
% =====================================================================
\section{Credibility tiers: definitions, the UMLS finding, and external alignment}
\label{sec:supp-credibility}

\subsection{What the tiers mean}
\label{sec:supp-credibility-defs}

The overall credibility rating is the model's only subjective core
dimension. Its four levels are defined by \emph{defeasibility}: what
would lead a curator to stop accepting an evidence item. They are not
defined by study design or by the certainty expressed by the authors:

\begin{description}
  \item[\dimval{VERY\_HIGH}] accepted by default even in the face of direct
    contradiction; reconsidered only in light of a stronger opposing
    argument;
  \item[\dimval{HIGH}] accepted unless directly contradicted by evidence of
    comparable strength; such a contradiction calls the item into question;
  \item[\dimval{MEDIUM}] not accepted on its own; accepted when corroborated
    by an independent source or an orthogonal type of evidence;
  \item[\dimval{LOW}] not relied upon; recorded pending independent
    replication.
\end{description}

The distinction between \dimval{HIGH} and \dimval{VERY\_HIGH} is
deliberate. Evidence of comparable strength is sufficient to call a
\dimval{HIGH} item into question, whereas a \dimval{VERY\_HIGH} item
remains the default position unless the opposing case is stronger.

This definition is a belief-revision policy. It therefore corresponds to
SEPIO's confidence attribute: the degree of belief that an asserted
proposition is true~\cite{sepio}. The credibility facets, including cohort
size, replication, multiple-testing control, ancestry control, and
ascertainment, inform that judgment but remain separate.

This separation is inspired by GRADE's general practice of assessing the
factors that affect certainty separately~\cite{grade2008}. GEM does not
adopt the GRADE framework or its rating rules, which were developed for
bodies of clinical evidence, particularly evidence from randomized and
observational studies. The GEM facets instead support credibility judgments
about individual evidence items from basic and pre-clinical research.

The scale intensifies at the top, as ACMG evidence strength
(\emph{supporting} to \emph{very strong})~\cite{acmg2015} and ClinGen
gene--disease validity (\emph{limited} to \emph{definitive})~\cite{clingen2018}
do. GRADE certainty instead extends downward to \emph{very low}, because it
rates a body of evidence by downgrading from a starting level determined by
study design. The GEM and GRADE scales therefore address related but
non-equivalent constructs; their alignment in \cref{tab:cred-alignment} is
deliberately non-positional.

\subsection{Why the tiers have no UMLS proxy}
\label{sec:supp-credibility-umls}

The four credibility tiers have no adequate UMLS proxy. This is a curation
conclusion, not merely a missing mapping: the tiers were tested in a
pre-registered sweep (\texttt{data/umls/sweeps/credibility.yaml}).

The adequacy criteria were fixed before the search. A candidate concept had
to (A) denote a degree of epistemic warrant rather than the magnitude of a
measured property; (B) represent a scale point rather than the scale
dimension itself; (C) belong to a set that is disjoint and ordered in the
same way as the GEM scale; and (D) carry that meaning in the relevant domain
rather than as a homonym.

The sweep used a local index of UMLS 2026AA and comprised 174 queries that
returned 497 distinct concepts. The results were replicated through the UTS
REST API, which returned 505 concepts. The search had three passes:
whole-string searches for intensified and degree qualifiers; searches for
the names of evidence-grading vocabularies; and searches for the same terms
within the UMLS semantic types most likely to contain scale values:
Qualitative Concept, Quantitative Concept, Idea or Concept, Intellectual
Product, and Finding. Each
retrieved concept was assigned the criterion it failed
(\cref{tab:cred-families}); the full generated table is available in the
results file.

\begin{table}[H]
  \centering
  \small
  \caption[Credibility sweep: rejected concept families]{Concepts retrieved by
  the credibility sweep (UMLS 2026AA), grouped by the adequacy criterion they
  failed. Criterion codes are defined in the text.}
  \label{tab:cred-families}
  \begin{tabularx}{\linewidth}{@{}>{\raggedright\arraybackslash}X r c >{\raggedright\arraybackslash}X@{}}
    \toprule
    Concept family & $n$ & Criterion failed & Examples \\
    \midrule
        “grade” (tumour, school, or product) & 202 & D & Astrocytoma, low grade (C1314694) \\
        LOINC/PROMIS self-efficacy items & 191 & D & Current level of confidence I can drive a car (C5142440) \\
        SNOMED degree qualifier values & 21 & A & Very high (C0442804); Extremely high (C5787630) \\
        clinical “evidence of” expressions (presence of a sign) & 20 & D & Blood alcohol level of less than 20 mg/100 ml (C0481372); Blood alcohol level of 20--39 mg/100 ml (C0481373) \\
        qualifiers denoting the degree of a measured property & 12 & A & Observation Value - High (C1561957); Message Waiting Priority - High (C1561958) \\
        scores, probabilities, risk categories & 8 & A & IPSS-R Risk Category Very High (C5202916); IPSS Risk Category High (C4522209) \\
        patient-perceived credibility of information & 7 & D & Establish teacher credibility, as appropriate (C0511083); credibility (C0870373) \\
        NCI/PDQ Level of Evidence I--IV & 6 & A & Level of Evidence (C0393009); Level of Evidence I (C2986612) \\
        metabolite-identification confidence levels & 6 & D & Level 1 Metabolite Identification Confidence (C5551417); Level 3 Metabolite Identification Confidence (C5551419) \\
        SNOMED diagnostic-certainty modality & 5 & C & Definite (C0439544); Probable diagnosis (C0332148) \\
        NCI confidence answer set & 4 & D & Very High Confidence (C4331500); High Confidence (C4330261) \\
        psychological self-confidence & 2 & D & Euphoric mood (C0235146); Confidence level (self-esteem) (C0518578) \\
        lexical collisions (“media”, “lower”) & 2 & D & Communications Media (C0009458); Culture Media (C0010454) \\
        degree of a measured property & 2 & A & A Medium Amount of Time (C4522282); A Medium Amount (C4522283) \\
        NCI adverse-event causality ladder & 2 & D & Definitely Related to Intervention (C1704787); Possibly Related to Intervention (C1705910) \\
        LOINC molecular-pathology “level of evidence” & 2 & D & Level of evidence:Find:Pt:Bld/Tiss:Nom (C4760312) \\
        diagnostic-certainty rating items & 2 & D & UHDRS 1999 Version - Diagnosis Confidence Level (C4744064); Level of Diagnostic Certainty (C5909371) \\
        TNM certainty factor & 2 & D & Degree of certainty of TNM classification (C0456878); Qualifier for TNM degree of certainty (C1302640) \\
        statistical confidence measures & 1 & A & Level of Confidence (statistical) (C5702556) \\
    \bottomrule
  \end{tabularx}
\end{table}

Four families merit explicit discussion because they are the closest
constructs retrieved from UMLS. First, the SNOMED qualifier values
\emph{High}~(C0205250, SNOMED 75540009), \emph{Very high}~(C0442804), and
\emph{Low}~(C0205251), together with related terms, occur under
\emph{Increased} and \emph{Decreased}. The source terminology therefore
treats them as degrees of a measured property, not as degrees of epistemic
warrant (A). Second, NCI Thesaurus \emph{Level of Evidence I--IV} is a
study-design hierarchy. Mapping GEM credibility to it would collapse the
model's separate \emph{method} and credibility dimensions (A, and C for the
scale). Third, the NCI \emph{Very High / High / Moderate / Low Confidence}
set has the same four-level, top-intensified form as the GEM scale, but it is
a \emph{Clinical or Research Assessment Answer}: a questionnaire response
that records a respondent's confidence in performing an activity (D).

The SNOMED diagnostic-certainty terms \emph{definite}, \emph{probable}, and
\emph{possible} are the closest epistemic candidates. They express the
likelihood that a proposition is true, but they form a three-level scale with
no intensified top category, apply in their source terminologies to
diagnostic assertions, and assess likelihood rather than defeasibility (C,
D). The same limitation applies to the TNM certainty factor, the LOINC
molecular-pathology \emph{level of evidence}, which is an
AMP/ASCO/CAP actionability tier for variant reports, and
metabolite-identification confidence levels. These are ordered epistemic scales, but
each is defined for a different object.

Lexical post-coordination does not supply the missing identifiers. Combining
a qualifier concept with an intensifier from the SPECIALIST Lexicon
(\texttt{LRMOD}) would add only a lexical modifier. The lexicon's
\emph{intensifier} class is a syntactic category that includes
\emph{very}, \emph{slightly}, \emph{about}, and \emph{ago}; it encodes
neither magnitude nor direction. Such an identifier therefore adds no
semantic content beyond the corresponding text string. This is a conceptual
argument about the role of identifiers, not an empirical finding.

\subsection{External alignment}
\label{sec:supp-credibility-align}

The following table gives a conceptual, non-positional alignment between the
GEM tiers and selected external scales. It does not define equivalences or a
common numerical conversion. In particular, the FHIR values are export codes:
\dimval{VERY\_HIGH} must be represented as \texttt{high}, and that export is
therefore lossy.

\begin{table}[H]
  \centering
  \small
  \caption[Credibility: external alignment]{Conceptual, non-positional
  alignment of the GEM credibility tiers with selected external scales.
  ``---'' indicates that no counterpart is proposed. The FHIR column gives the
  lossy export used when a \texttt{certainty-rating} code is
  required~\cite{fhir_evidence}. CIO = Confidence Information
  Ontology~\cite{cio2015}.}
  \label{tab:cred-alignment}
  \begin{tabularx}{\linewidth}{@{}l >{\raggedright\arraybackslash}X >{\raggedright\arraybackslash}X >{\raggedright\arraybackslash}X >{\raggedright\arraybackslash}X >{\raggedright\arraybackslash}X@{}}
    \toprule
    GEM & GRADE certainty~\cite{grade2008} & FHIR \texttt{certainty-rating} & CIO confidence level~\cite{cio2015} & ACMG evidence strength~\cite{acmg2015} & ClinGen validity~\cite{clingen2018} \\
    \midrule
    \dimval{VERY\_HIGH} & --- (GRADE caps at \emph{high}) & \texttt{high} (lossy export) & high (partial) & very
    strong & definitive \\
    \dimval{HIGH}       & high     & \texttt{high}     & high   & strong     & strong \\
    \dimval{MEDIUM}     & moderate & \texttt{moderate} & medium & moderate   & moderate \\
    \dimval{LOW}        & low      & \texttt{low}      & low (``absence of trust'', closer to none) & supporting &
    limited \\
    ---                 & very low & \texttt{very-low} & ---    & ---        & --- \\
    \bottomrule
  \end{tabularx}
\end{table}

These alignments are part of the GEM model, not of the UMLS crosswalk. The
crosswalk records UMLS mappings, or a considered conclusion that no adequate
mapping exists.
